\documentclass[11pt,a4paper]{report}
\usepackage{fontspec}\usepackage[spanish,es-nodecimaldot]{babel}\usepackage{geometry}\usepackage{graphicx}\usepackage{float}\usepackage{xcolor}\usepackage{hyperref}\usepackage{listings}\usepackage{longtable}\usepackage{booktabs}\usepackage{tikz}\usepackage{enumitem}\usepackage{fancyhdr}\usepackage{dirtree}\usepackage{caption}
\geometry{margin=2.4cm}\setmainfont{TeX Gyre Pagella}\setsansfont{TeX Gyre Heros}\definecolor{navy}{HTML}{08243A}\definecolor{cyan}{HTML}{129E98}\definecolor{codebg}{HTML}{F2F5F7}
\hypersetup{colorlinks=true,linkcolor=navy,urlcolor=cyan}\pagestyle{fancy}\fancyhf{}\lhead{TW MaintFlow}\rhead{Guía de defensa}\cfoot{\thepage}
\lstset{basicstyle=\ttfamily\small,backgroundcolor=\color{codebg},frame=single,breaklines=true,showstringspaces=false,keywordstyle=\color{navy}\bfseries,commentstyle=\color{cyan}}
\newcommand{\figura}[3]{\begin{figure}[H]\centering\includegraphics[width=#1\textwidth]{../#2}\caption{#3}\label{fig:#2}\end{figure}}
\begin{document}
\begin{titlepage}\pagecolor{navy}\color{white}\centering\vspace*{1.2cm}{\Large UNIVERSIDAD NACIONAL DEL CALLAO\par}\vspace{.4cm}{\large Facultad de Ingeniería Industrial y de Sistemas\par Escuela Profesional de Ingeniería de Sistemas\par}\vfill{\Huge\bfseries TW MaintFlow\par}\vspace{.5cm}{\Large Sistema web MVC para gestión de mantenimiento y órdenes de servicio\par}\vfill\begin{tabular}{rl}Curso:&Programación Web I - ISEE405\\Docente:&Mg. Junior Grados Espinoza\\Estudiante:&Jhon Gesell Villanueva Portella\\Código:&2515250011\\Material compartido con:&Maricielo\\Año:&2026\end{tabular}\vfill{\large Callao, Perú\par}\end{titlepage}\nopagecolor\color{black}
\tableofcontents\listoffigures\listoftables
\chapter{Resumen del proyecto}
TW MaintFlow resuelve la dispersión de información en organizaciones que mantienen equipos industriales. Centraliza clientes, sedes, equipos, técnicos, órdenes, repuestos, usuarios y reportes. Sus usuarios son administración, recepción y personal técnico. El alcance académico usa datos sintéticos y demuestra trazabilidad, control de inventario y separación MVC con HTML5, CSS3, JavaScript, PHP, MariaDB y XAMPP.
\chapter{Tecnologías usadas}
HTML5 estructura formularios y tablas; CSS3 aporta diseño responsive; JavaScript confirma desactivaciones y valida cantidades; PHP procesa rutas, sesiones y reglas; PDO ejecuta sentencias preparadas; MariaDB persiste 12 entidades; Apache sirve la aplicación; XAMPP reúne el entorno y phpMyAdmin administra la base. \texttt{password\_hash} genera hashes y \texttt{password\_verify} los comprueba sin almacenar contraseñas en texto plano.
\chapter{Teoría mínima de MVC}
MVC divide responsabilidades. El Modelo conoce datos y consultas; la Vista presenta HTML; el Controlador coordina petición, permisos, modelo y respuesta. Frente a PHP estructurado mezclado, mejora mantenibilidad, reutilización y pruebas: cambiar el diseño no obliga a reescribir SQL y cambiar una consulta no exige modificar la vista.
\figura{.95}{figuras/03_componentes_mvc.png}{Responsabilidades de los componentes MVC. Fuente: elaboración propia.}
\chapter{Arquitectura real del proyecto}
\begin{lstlisting}
tw_maintflow/
|-- app/{config,core,controllers,models,views}
|-- public/
|-- database/
|-- storage/
`-- documentacion_defensa/
\end{lstlisting}
\texttt{app/config} contiene configuración y PDO; \texttt{core} contiene App, Router, Controller, Model, Session y Auth; \texttt{controllers} coordina casos de uso; \texttt{models} contiene persistencia; \texttt{views} genera la interfaz; \texttt{public} es el punto expuesto; \texttt{database} instala y consulta MariaDB.
\figura{.95}{figuras/01_arquitectura_general.png}{Arquitectura general. Fuente: elaboración propia.}
\chapter{Flujo de una petición}
El navegador solicita una ruta; Apache aplica \texttt{.htaccess}; \texttt{public/index.php} carga configuración y autoload; Router separa controlador y acción; el controlador valida sesión/rol; instancia un modelo; PDO consulta MariaDB; el controlador entrega datos a la vista; la vista genera HTML y el navegador lo representa.
\figura{.95}{figuras/02_flujo_mvc.png}{Ejemplo real de la ruta clientes. Fuente: elaboración propia.}
\chapter{Base de datos}
\begin{longtable}{p{3.2cm}p{2.2cm}p{4.2cm}p{4cm}}\toprule Tabla&PK&FK principales&Propósito\\\midrule\endhead
roles&id\_rol&--&Perfiles de autorización\\usuarios&id\_usuario&id\_rol&Credenciales y datos personales\\clientes&id\_cliente&--&Empresas atendidas\\sedes&id\_sede&id\_cliente&Locales de clientes\\categorias\_equipo&id\_categoria&--&Clasificación de activos\\equipos&id\_equipo&id\_sede, id\_categoria&Activos mantenibles\\tecnicos&id\_tecnico&id\_usuario&Perfil técnico\\estados\_orden&id\_estado&--&Catálogo de estados\\ordenes\_servicio&id\_orden&id\_equipo, id\_tecnico, id\_estado&Cabecera del servicio\\detalle\_orden&id\_detalle&id\_orden, id\_repuesto&Actividad y consumo\\repuestos&id\_repuesto&--&Inventario\\movimientos\_repuesto&id\_movimiento&id\_repuesto, id\_usuario&Kardex\\\bottomrule\end{longtable}
Ejemplos: una sede ``Frigorífico'', el equipo EQ-009 y una orden correctiva. Todos son sintéticos.
\chapter{Relaciones principales}
Un rol tiene usuarios; un cliente tiene sedes; una sede tiene equipos; una categoría clasifica equipos; un usuario puede originar un técnico; técnicos, equipos y estados se relacionan con órdenes; una orden tiene detalles; repuestos aparecen en detalles y movimientos; el usuario identifica al responsable del movimiento.
\figura{.98}{figuras/04_modelo_relacional.png}{Modelo relacional simplificado. Fuente: elaboración propia.}
\chapter{Módulos funcionales}
El inicio presenta el producto; login autentica; dashboard resume indicadores; clientes, sedes, equipos, técnicos y repuestos usan CRUD; órdenes registran y actualizan trabajo; usuarios y reportes son administrativos; logout destruye la sesión. Los errores 403, 404 y 500 distinguen falta de permiso, ruta inexistente y fallo interno.
\chapter{Roles y permisos}
Administrador: CRUD general, usuarios y reportes. Recepción: clientes, sedes, equipos y creación/consulta de órdenes. Técnico: consulta órdenes y registra estado, diagnóstico, solución y observaciones. La interfaz oculta opciones y Auth vuelve a validar en servidor.
\figura{.98}{figuras/07_roles_permisos.png}{Matriz de permisos implementada. Fuente: elaboración propia.}
\chapter{CRUD}
CRUD significa Create, Read, Update y Delete. El proyecto prefiere desactivación lógica. En clientes, \texttt{/clientes} ejecuta \texttt{ClienteController::index}; \texttt{CrudModel::all} lee; la vista lista; \texttt{/clientes/form} presenta formulario; POST a \texttt{/clientes/save} filtra campos permitidos y ejecuta INSERT o UPDATE preparado; finalmente redirige.
\chapter{Autenticación y sesiones}
\texttt{Session::start} inicia sesión. AuthController recibe usuario y contraseña, Usuario busca un usuario activo y \texttt{password\_verify} compara el hash. Al éxito se regenera el ID y se guarda id, nombre y rol. \texttt{requireLogin} protege rutas y \texttt{requireRole} responde 403. Logout vacía datos, elimina cookie y destruye sesión.
\figura{.95}{figuras/05_flujo_login.png}{Flujo de autenticación. Fuente: elaboración propia.}
\chapter{Órdenes de servicio}
Recepción selecciona equipo, técnico, estado, tipo, prioridad, fecha y falla. El técnico actualiza diagnóstico, solución y estado. La orden posee detalles con actividad, repuesto, cantidad, precio y subtotal generado. El costo total combina mano de obra y detalles; al estado 5 se registra fecha de cierre.
\figura{.95}{figuras/06_flujo_orden_servicio.png}{Ciclo de una orden. Fuente: elaboración propia.}
\chapter{Trigger, procedimiento y función}
\begin{table}[H]\centering\begin{tabular}{llll}\toprule Objeto&Implementado&Propósito&Archivo\\\midrule tr\_detalle\_stock&Sí&Descontar stock y recalcular total&tw\_maintflow.sql\\sp\_cerrar\_orden&Sí&Cerrar orden y fechar&tw\_maintflow.sql\\fn\_total\_orden&Sí&Calcular total por id&tw\_maintflow.sql\\\bottomrule\end{tabular}\caption{Objetos SQL comprobados en el script real.}\end{table}
El trigger es AFTER INSERT sobre \texttt{detalle\_orden}; si hay repuesto y cantidad positiva, resta stock y recalcula costo. Su aporte es consistencia automática; su riesgo es permitir stock negativo porque no existe validación previa. La prueba insertó una unidad en transacción: 18 pasó a 17 y luego se hizo rollback.
\begin{lstlisting}[language=SQL]
CREATE TRIGGER tr_detalle_stock AFTER INSERT ON detalle_orden
FOR EACH ROW BEGIN
 IF NEW.id_repuesto IS NOT NULL AND NEW.cantidad>0 THEN
  UPDATE repuestos SET stock_actual=stock_actual-NEW.cantidad
  WHERE id_repuesto=NEW.id_repuesto;
 END IF;
END
\end{lstlisting}
\figura{.95}{figuras/08_trigger_stock.png}{Flujo del trigger. Fuente: elaboración propia.}
\chapter{Consultas SQL importantes}
\texttt{consultas\_demo.sql} contiene INNER JOIN para unir órdenes, equipos y estados; LEFT JOIN para conservar técnicos sin órdenes; COUNT y GROUP BY para totales; SUM para consumo/costo; AVG para costo promedio; HAVING filtra grupos; ORDER BY ordena costos; BETWEEN filtra fechas; y \texttt{stock\_actual <= stock\_minimo} detecta alertas.
\begin{lstlisting}[language=SQL]
SELECT es.nombre_estado, COUNT(*) total
FROM ordenes_servicio o
JOIN estados_orden es ON es.id_estado=o.id_estado
GROUP BY es.id_estado;
\end{lstlisting}
\chapter{Seguridad básica}
Se emplean hashes, prepared statements, \texttt{htmlspecialchars}, validación de sesión y rol, IDs con \texttt{FILTER\_VALIDATE\_INT}, lista blanca de tablas/campos, confirmación JavaScript y datos sintéticos. Mejora pendiente: token CSRF y reglas de validación más profundas.
\chapter{Fragmentos de código explicados}
\section*{Punto de entrada: public/index.php}
\begin{lstlisting}[language=PHP]
require_once dirname(__DIR__).'/app/config/config.php';
require_once APP_ROOT.'/app/config/database.php';
spl_autoload_register(function(string $class): void { /* carga clases */ });
Session::start();
header('Content-Type: text/html; charset=UTF-8');
(new App())->run();
\end{lstlisting}
Carga configuración, registra clases, abre sesión, declara UTF-8 y delega a App.
\section*{Router y controlador base}
\begin{lstlisting}[language=PHP]
$parts = array_values(array_filter(explode('/', $path)));
$class = $map[$segment].'Controller';
$action = $parts[1] ?? 'index';
(new $class)->$action();
\end{lstlisting}
El mapa impide instanciar clases arbitrarias. Controller extrae datos, incluye header/vista/footer y ofrece redirección.
\section*{Modelo, Auth y conexión}
\begin{lstlisting}[language=PHP]
abstract class Model {
 protected PDO $db;
 public function __construct(){ $this->db=Database::connection(); }
}
// DSN real: mysql:host=localhost;dbname=bd_tw_maintflow;charset=utf8mb4
\end{lstlisting}
Todos los modelos heredan la misma conexión. Auth comprueba sesión y roles antes de ejecutar casos de uso.
\section*{Clientes y órdenes}
ClienteController configura tabla, PK, carpeta y campos; hereda el flujo de CrudController. Cliente hereda CrudModel. La vista compartida escapa cada valor. OrdenServicio incluye JOIN, opciones, creación y actualización técnica con sentencias preparadas. Los doce fragmentos ampliados están en \texttt{codigo/fragmentos\_explicados.txt}.
\chapter{Diagramas}
Los ocho diagramas anteriores muestran arquitectura, flujo MVC, componentes, relaciones, autenticación, ciclo de orden, roles y trigger. Se generaron como PNG legibles y se conservan en \texttt{documentacion\_defensa/figuras}.
\chapter{Capturas reales del sistema}
\foreach \n/\t in {01_inicio/Inicio público,02_login/Login,03_dashboard_admin/Dashboard administrativo,04_clientes/Clientes,05_sedes_utf8_corregido/Sedes con UTF-8,06_equipos/Equipos,07_tecnicos/Técnicos,08_ordenes/Órdenes,09_detalle_orden/Detalle de orden,10_repuestos/Repuestos,11_usuarios/Usuarios,12_reportes/Reportes,13_error_403/Error 403,14_error_404/Error 404,15_phpmyadmin_tablas/Tablas en phpMyAdmin}{\figura{.92}{capturas/\n.png}{\t. Fuente: captura real de TW MaintFlow.}}
\chapter{Preguntas posibles del docente}
\begin{enumerate}[leftmargin=*]
\item ¿Qué es MVC? Separación entre datos, presentación y coordinación.
\item ¿Modelo vs. controlador? El modelo persiste; el controlador dirige el caso de uso.
\item ¿Por qué no hay SQL en vistas? Para desacoplar presentación y datos.
\item ¿Qué hace Router? Traduce URL a controlador y acción permitidos.
\item ¿Qué es sesión? Estado de usuario conservado en el servidor.
\item ¿Cómo se validan roles? Con \texttt{Auth::requireRole}.
\item ¿Qué hace password\_verify? Compara contraseña con hash.
\item ¿Qué es PK? Identificador único de una fila.
\item ¿Qué es FK? Restricción que referencia otra tabla.
\item ¿Por qué 12 tablas? Para normalizar las entidades del dominio.
\item ¿Cliente-sede? Uno a muchos.
\item ¿Sede-equipo? Uno a muchos.
\item ¿Equipo-orden? Uno a muchos.
\item ¿Orden-detalle? Uno a muchos.
\item ¿Cómo funciona el trigger? Después del detalle resta stock y recalcula costo.
\item ¿Riesgo del trigger? No impide stock negativo.
\item ¿Procedimiento almacenado? Rutina SQL invocable; aquí cierra una orden.
\item ¿Función SQL? Rutina que retorna valor; aquí retorna total.
\item ¿Por qué son opcionales? MVC funciona sin objetos almacenados.
\item ¿Prepared statement? SQL compilado con parámetros separados.
\item ¿Cómo evita inyección? No concatena valores dentro de SQL.
\item ¿Ruta inexistente? Router responde 404.
\item ¿403 vs. 404? Prohibido vs. no encontrado.
\item ¿Cómo se descuenta stock? Trigger AFTER INSERT.
\item ¿Cómo se calcula costo? Mano de obra más suma de subtotales.
\item ¿Por qué utf8mb4? Representa Unicode completo y español correctamente.
\item ¿GET vs. POST? GET consulta; POST envía cambios.
\item ¿Qué es CRUD? Crear, leer, actualizar y eliminar/desactivar.
\item ¿GROUP BY? Agrupa filas para agregados.
\item ¿WHERE vs. HAVING? Filtra filas antes y grupos después.
\item ¿INNER JOIN? Devuelve coincidencias.
\item ¿LEFT JOIN? Conserva todas las filas izquierdas.
\item ¿Qué hace htmlspecialchars? Escapa HTML para reducir XSS.
\item ¿Qué hace session\_regenerate\_id? Cambia ID tras login.
\item ¿Qué es PDO? API uniforme de acceso a bases.
\item ¿Qué hace .htaccess? Reescribe rutas hacia index.php.
\item ¿Fallback de rutas? \texttt{index.php?url=clientes/index}.
\item ¿Qué patrón usa public/index.php? Front Controller.
\item ¿Cómo se cierra sesión? Vacía sesión, cookie y destruye estado.
\item ¿Mejora prioritaria? CSRF, validación y control de stock previo.
\end{enumerate}
\chapter{Guion de exposición de siete minutos}
\begin{enumerate}\item[0:00] Problema: información dispersa de mantenimiento.\item[0:40] Solución y módulos.\item[1:20] Tecnologías XAMPP/PHP/MariaDB.\item[2:00] MVC y flujo de una ruta.\item[3:00] Doce tablas y relaciones.\item[3:50] Login y tres roles.\item[4:40] Ciclo de orden.\item[5:30] Trigger, procedimiento y función.\item[6:10] Dashboard/reportes.\item[6:40] Cumplimiento, seguridad y mejoras.\end{enumerate}
\chapter{Práctica guiada}
Jhon y Maricielo deben localizar ClienteController, Cliente/CrudModel y la vista compartida; seguir \texttt{/clientes}; explicar CRUD y una FK; ejecutar un JOIN; probar cada login; provocar 403 con técnico y 404 con una ruta inventada; crear un cliente y una orden; insertar un detalle en transacción y comprobar el descuento de stock.
\chapter{Glosario}
\begin{description}[style=nextline]\item[MVC] Modelo-Vista-Controlador.\item[CRUD] Operaciones de persistencia.\item[Router] Resolutor de rutas.\item[Controller] Coordinador.\item[Model] Acceso a datos.\item[View] Presentación.\item[Session] Estado de navegación.\item[Authentication] Verificación de identidad.\item[Authorization] Verificación de permisos.\item[Trigger] Acción automática SQL.\item[Procedure] Rutina SQL sin retorno obligatorio.\item[Function] Rutina SQL con retorno.\item[Primary key] Clave primaria.\item[Foreign key] Clave foránea.\item[JOIN] Unión de tablas.\item[Hash] Representación irreversible.\item[SQL Injection] Inyección de código SQL.\item[UTF-8] Codificación Unicode.\item[utf8mb4] Juego completo UTF-8 en MariaDB.\end{description}
\chapter{Conclusiones}
El proyecto cumple una arquitectura MVC educativa, 12 tablas relacionadas y ejecución en XAMPP. Controladores, modelos y vistas están separados; sesiones y roles protegen rutas; UTF-8 se aplica de extremo a extremo. El trabajo integra conceptos del curso y deja como mejoras CSRF, validación avanzada, auditoría, control preventivo de stock y pruebas automatizadas.
\end{document}
